% =============================================================================
%  APÊNDICE C: CÓDIGOS COMPLEMENTARES
% =============================================================================
\chapter{Códigos Complementares}
\label{apend:codigos}

Este apêndice reúne implementações completas de exemplos e componentes
mencionados nos capítulos. Os códigos são extraídos do código-fonte real
do \opencodeeco{} e podem ser executados diretamente.

% =============================================================================
%  SEÇÃO C.1: SCANNER PERSONALIZADO
% =============================================================================
\section{Exemplo de Implementação de um Scanner}
\label{sec:cod-scanner}

O código abaixo implementa um scanner epistemológico simplificado que
identifica lacunas de capacidade no ecossistema. Corresponde ao padrão
utilizado pelo Noological Scanner (SPEC-028).

\begin{lstlisting}[caption={Scanner epistemologico simplificado},
                   label={cod:scanner}, style=pythonstyle]
"""
noological_scanner_demo.py — Scanner epistemologico simplificado.

Demonstra o padrao de implementacao de um scanner no OpenCode Ecosystem.
"""

from dataclasses import dataclass, field
from typing import List, Optional


@dataclass
class Capability:
    """Representa uma capacidade ou componente do ecossistema."""
    name: str
    category: str          # skill, mcp, agent, plugin
    description: str
    dependencies: List[str] = field(default_factory=list)
    implemented: bool = True


@dataclass
class CapabilityGap:
    """Representa uma lacuna de capacidade identificada."""
    name: str
    rationale: str
    priority: int          # 1 (critica) a 5 (desejavel)
    suggested_category: str
    dependencies: List[str] = field(default_factory=list)


class NoologicalScanner:
    """
    Scanner que identifica 'o que nao existe' no ecossistema.
    """

    def __init__(self, registry: List[Capability]):
        self.registry = {c.name: c for c in registry}
        self.gaps: List[CapabilityGap] = []

    def scan(self, required: List[str]) -> List[CapabilityGap]:
        """
        Varre a lista de capacidades requeridas e identifica quais
        nao estao implementadas.

        Args:
            required: Lista de nomes de capacidades requeridas.

        Returns:
            Lista de lacunas identificadas.
        """
        self.gaps.clear()
        for name in required:
            if name not in self.registry:
                gap = CapabilityGap(
                    name=name,
                    rationale=f"Capacidade '{name}' requerida mas "
                              f"nao implementada",
                    priority=1,
                    suggested_category="skill",
                    dependencies=[]
                )
                self.gaps.append(gap)
        return self.gaps

    def suggest_implementation(self, gap: CapabilityGap) -> str:
        """Gera esboco de implementacao para a lacuna."""
        return (
            f"# {gap.name}\n"
            f"# Prioridade: {gap.priority}\n"
            f"# Dependencias: {', '.join(gap.dependencies) or 'Nenhuma'}\n\n"
            "class NovaSkill:\n"
            f'    """Skill gerada pelo scanner: {gap.name}."""\n\n'
            "    def execute(self, context):\n"
            '        """Implementacao pendente."""\n'
            "        raise NotImplementedError\n"
        )


# Exemplo de uso
if __name__ == "__main__":
    registry = [
        Capability("web-search", "skill",
                   "Busca na web via DuckDuckGo"),
        Capability("code-runner", "skill",
                   "Execucao de codigo Python"),
    ]
    scanner = NoologicalScanner(registry)
    gaps = scanner.scan([
        "web-search", "code-runner",
        "academic-search", "pdf-extractor"
    ])

    for g in gaps:
        print(f"[LACUNA] {g.name} (prioridade {g.priority})")
        print(scanner.suggest_implementation(g))
        print("---")
\end{lstlisting}

% =============================================================================
%  SEÇÃO C.2: CONFIGURAÇÃO DO TRUST ENGINE
% =============================================================================
\section{Exemplo de Configuração do Trust Engine}
\label{sec:cod-trust}

A configuração abaixo ilustra a inicialização e uso do Trust Engine
(SPEC-038) com todos os seus componentes.

\begin{lstlisting}[caption={Configuracao do Trust Engine},
                   label={cod:trust}, style=pythonstyle]
"""
trust_engine_demo.py — Configuracao do Trust Engine.

Demonstra TrustScorer (blend 70/30), BehavioralGate,
NaturalForgetting e OutcomeTracker.
"""

from enum import Enum, auto
from dataclasses import dataclass, field
from typing import List, Dict


class ActionCategory(Enum):
    """Categorias de acao para classificacao comportamental."""
    READ = auto()
    WRITE = auto()
    EXECUTE = auto()
    NETWORK = auto()
    GOVERNANCE = auto()


class SafetyLevel(Enum):
    """Niveis de seguranca definidos pelo BehavioralGate."""
    SAFE = "safe"
    MODERATE = "moderate"
    RISKY = "risky"
    BLOCKED = "blocked"


@dataclass
class Action:
    """Representa uma acao executada por um agente."""
    agent_id: str
    category: ActionCategory
    description: str
    timestamp: float = 0.0


@dataclass
class TrustScore:
    """Pontuacao de confianca de um agente."""
    direct_score: float = 0.85       # 70% do peso
    network_score: float = 0.80      # 30% do peso
    history: List[float] = field(default_factory=list)

    @property
    def blended(self) -> float:
        """Blend 70/30 entre score direto e de rede."""
        return 0.7 * self.direct_score + 0.3 * self.network_score


class BehavioralGate:
    """Barreira preventiva que classifica acoes antes da execucao."""

    RULES: Dict[ActionCategory, SafetyLevel] = {
        ActionCategory.READ: SafetyLevel.SAFE,
        ActionCategory.WRITE: SafetyLevel.MODERATE,
        ActionCategory.EXECUTE: SafetyLevel.RISKY,
        ActionCategory.NETWORK: SafetyLevel.MODERATE,
        ActionCategory.GOVERNANCE: SafetyLevel.SAFE,
    }

    def check(self, action: Action) -> SafetyLevel:
        """
        Verifica a seguranca de uma acao.

        Args:
            action: Acao a ser verificada.

        Returns:
            Nivel de seguranca classificado.
        """
        base = self.RULES.get(action.category, SafetyLevel.MODERATE)
        return base


class TrustEngine:
    """
    Motor de confianca completo: scorer + gate + forgetting + tracker.
    """

    def __init__(self):
        self.scores: Dict[str, TrustScore] = {}
        self.gate = BehavioralGate()
        self.history: List[dict] = []

    def register_agent(self, agent_id: str) -> None:
        """Registra um novo agente com score inicial."""
        self.scores[agent_id] = TrustScore()

    def evaluate(self, agent_id: str,
                 action: Action) -> tuple[SafetyLevel, float]:
        """
        Avalia uma acao: classifica seguranca e retorna confianca.

        Args:
            agent_id: Identificador do agente.
            action: Acao a ser avaliada.

        Returns:
            Tupla (nivel de seguranca, confianca blendada).
        """
        safety = self.gate.check(action)
        score = self.scores.get(agent_id, TrustScore())
        self.history.append({
            "agent": agent_id,
            "action": action.description,
            "safety": safety.value,
            "trust": score.blended,
        })
        return safety, score.blended
\end{lstlisting}

% =============================================================================
%  SEÇÃO C.3: AGENTE PERSONALIZADO
% =============================================================================
\section{Exemplo de Agente Personalizado}
\label{sec:cod-agente}

O código seguinte demonstra a criação de um agente cognitivo personalizado
no \opencodeeco{}, incluindo definição de skill, MCP e ciclo percepção-ação.

\begin{lstlisting}[caption={Agente personalizado},
                   label={cod:agente}, style=pythonstyle]
"""
custom_agent_demo.py — Agente personalizado no OpenCode Ecosystem.

Demonstra criacao de agente com skill, MCP e ciclo percepcao-acao.
"""

from dataclasses import dataclass
from typing import Optional


@dataclass
class Context:
    """Contexto de execucao do agente."""
    user_input: str
    history: list = None
    trust_score: float = 0.85


class Skill:
    """Habilidade executavel por um agente."""

    def __init__(self, name: str, version: str = "1.0.0"):
        self.name = name
        self.version = version

    def execute(self, context: Context) -> str:
        """Executa a skill no contexto fornecido."""
        raise NotImplementedError


class SearchSkill(Skill):
    """Skill de busca academica."""

    def __init__(self):
        super().__init__("academic-search", "2.1.0")

    def execute(self, context: Context) -> str:
        query = context.user_input
        return f"[Resultados da busca: '{query}' — 5 papers encontrados]"


class CustomAgent:
    """Agente cognitivo personalizado."""

    def __init__(self, agent_id: str, skills: list[Skill]):
        self.agent_id = agent_id
        self.skills = {s.name: s for s in skills}

    def perceive(self, context: Context) -> str:
        """Fase de percepcao: interpreta a entrada do usuario."""
        return f"Processando: {context.user_input}"

    def reason(self, perception: str) -> str:
        """Fase de raciocinio: decide qual skill usar."""
        if "buscar" in perception.lower():
            return "academic-search"
        return "default"

    def act(self, skill_name: str, context: Context) -> str:
        """Fase de acao: executa a skill selecionada."""
        skill = self.skills.get(skill_name)
        if not skill:
            return f"[Agente {self.agent_id}: skill '{skill_name}' nao encontrada]"
        return skill.execute(context)

    def run(self, user_input: str) -> str:
        """Ciclo completo percepcao--raciocinio--acao."""
        ctx = Context(user_input=user_input)
        perception = self.perceive(ctx)
        chosen = self.reason(perception)
        result = self.act(chosen, ctx)
        return result


# Exemplo de uso
if __name__ == "__main__":
    agent = CustomAgent("agente-demo", [SearchSkill()])
    output = agent.run("buscar algoritmos de confianca computacional")
    print(output)
\end{lstlisting}

% =============================================================================
%  SEÇÃO C.4: BENCHMARK CORA-EVAL
% =============================================================================
\section{Script de Benchmark CORA-Eval}
\label{sec:cod-cora}

O script abaixo implementa o rastreador evolutivo do benchmark CORA-Eval
(Capítulo~7), com 150 tarefas em 10 dimensões e 4 níveis.

\begin{lstlisting}[caption={CORA-Eval benchmark tracker},
                   label={cod:cora}, style=pythonstyle]
"""
cora_benchmark_tracker.py — Rastreador do benchmark CORA-Eval.

Avalia ecossistemas cognitivos em 150 tarefas distribuidas em
10 dimensoes × 4 niveis (Basico, Intermediario, Avancado, Pesquisa).
"""

from dataclasses import dataclass, field
from typing import Dict, List
import json


@dataclass
class CORATask:
    """Tarefa do benchmark CORA-Eval."""
    id: str                   # e.g., "LOG-01"
    dimension: str            # e.g., "Logica"
    level: int                # 1-4 (Basico a Pesquisa)
    description: str
    weight: float = 1.0


@dataclass
class CORAResult:
    """Resultado de avaliacao de uma tarefa."""
    task_id: str
    score: float              # 0.0 a 1.0
    passed: bool
    details: str = ""


class CORAEvalTracker:
    """Rastreador de benchmark com persistencia JSON."""

    CORA_DIMENSIONS = [
        "Logica", "Probabilidade", "Algebra", "Calculo",
        "Grafos", "Aprendizado", "Raciocinio", "Agentes",
        "Etica", "Inovacao"
    ]

    def __init__(self):
        self.tasks: Dict[str, CORATask] = {}
        self.results: Dict[str, CORAResult] = {}
        self._seed_tasks()

    def _seed_tasks(self) -> None:
        """Popula as 150 tarefas do benchmark."""
        for dim_idx, dim in enumerate(self.CORA_DIMENSIONS):
            for level in range(1, 5):
                count = 4 if level < 4 else 3
                for i in range(count):
                    tid = f"{dim[:3].upper()}-{level}{chr(65+i)}"
                    self.tasks[tid] = CORATask(
                        id=tid,
                        dimension=dim,
                        level=level,
                        description=f"Tarefa {dim} nivel {level} #{i+1}",
                        weight=1.0 + (level - 1) * 0.25,
                    )

    def record(self, task_id: str, score: float,
               details: str = "") -> CORAResult:
        """Registra o resultado de uma tarefa."""
        result = CORAResult(
            task_id=task_id,
            score=score,
            passed=score >= 0.7,
            details=details,
        )
        self.results[task_id] = result
        return result

    @property
    def cora_score(self) -> float:
        """CORA-Score: media ponderada por nivel."""
        if not self.results:
            return 0.0
        total_weight = sum(
            self.tasks[t].weight
            for t in self.results
            if t in self.tasks
        )
        if total_weight == 0:
            return 0.0
        weighted = sum(
            r.score * self.tasks[r.task_id].weight
            for r in self.results.values()
            if r.task_id in self.tasks
        )
        return weighted / total_weight

    def save(self, path: str = "cora_results.json") -> None:
        """Persiste os resultados em JSON."""
        data = {
            "cora_score": self.cora_score,
            "results": [
                {"task_id": r.task_id, "score": r.score,
                 "passed": r.passed}
                for r in self.results.values()
            ]
        }
        with open(path, "w", encoding="utf-8") as f:
            json.dump(data, f, indent=2, ensure_ascii=False)


# Exemplo de uso
if __name__ == "__main__":
    tracker = CORAEvalTracker()
    tracker.record("LOG-1A", 0.95, "Resolvido com Z3")
    tracker.record("PRO-2B", 0.80, "Inferencia bayesiana correta")
    print(f"CORA-Score: {tracker.cora_score:.3f}")
    tracker.save()
\end{lstlisting}

% =============================================================================
%  SEÇÃO C.5: COMANDOS MAKEFILE
% =============================================================================
\section{Comandos Makefile}
\label{sec:cod-makefile}

O \opencodeeco{} utiliza um \texttt{Makefile} para automatizar tarefas
recorrentes. Abaixo os principais comandos.

\begin{lstlisting}[caption={Comandos Makefile do OpenCode Ecosystem},
                   label={cod:makefile}, style=bashstyle]
# ============================================================
# Makefile - OpenCode Ecosystem
# Comandos: desenvolvimento, testes e compilacao
# ============================================================

# ---- COMPILACAO DO LIVRO ----
make livro

# Compilar apenas capitulo especifico
make capitulo CAP=08-capitulo1-mat-est

# Limpar arquivos auxiliares
make clean

# ---- TESTES ----
# Executar suite completa de testes (312 CTs)
make test

# Executar testes de um SPEC especifico
make test-spec SPEC=SPEC-038

# Executar testes com relatorio de cobertura
make test-cov

# ---- ECOSSISTEMA ----
# Sincronizar ecossistema (descobrir e instalar skills)
make sync

# Executar scanner pipeline completo
make scan

# Atualizar componentes
make update

# ---- CORA-EVAL ----
# Executar benchmark CORA-Eval completo
make cora-eval

# Executar benchmark por dimensao
make cora-dim DIM=Logica

# Visualizar resultados do benchmark
make cora-report

# ---- DOCUMENTACAO ----
# Gerar documentacao do ecossistema
make docs

# Verificar SPECs e ADRs
make validate-specs

# Atualizar graficos TikZ
make tikz

# ---- LIMPEZA ----
# Limpar completamente (incluindo PDF)
make distclean

# Limpar resultados de benchmark
make clean-bench
\end{lstlisting}
